\section{总结与改进方向}

本实验完成了一个较完整的鸡尾酒会语音分离系统。数据层面，实验从 LibriSpeech 中整理真实说话人语音，构造 Easy 和 Hard 两组双说话人混合数据，并在 Hard 实验中加入 ESC-50 真实环境噪声。模型层面，实验采用 STFT 特征、BLSTM 时序建模、软掩码预测和 PIT 损失，实现了从混合语音到两路分离语音的端到端处理流程。评估层面，实验使用 SI-SDR、Delta SDR 和最优排列匹配衡量分离效果。

实验结果表明，模型在 Easy 验证集上取得了明显改善，完整验证集平均 \texttt{Delta\_SDR=6.966 dB}，精选展示样例的 \texttt{SDR\_out} 可达到 13 dB 以上，主观听感较清晰。Hard 实验由于引入更复杂说话人池、等能量混合和真实环境噪声，完整验证集平均 \texttt{Delta\_SDR=2.549 dB}，低于 Easy，但最佳展示样例仍能达到 10 dB 以上的 \texttt{SDR\_out}。这说明模型具备基础分离能力，但在噪声和强重叠条件下仍有提升空间。

同时，训练集与验证集对比显示，Easy 和 Hard 模型的验证效果都没有训练集效果好，说明当前模型的泛化能力有限。这一点与个人设备、训练时间和数据规模限制有关：在有限样本上训练 200 个 epoch 能够让模型较好拟合训练集，但在未见过的说话人与噪声组合上仍容易下降。不过，从 Oracle 理想掩码结果以及部分验证样例的高质量分离效果可以看出，数据和程序流程本身具有较高理论上限，后续通过扩大数据、延长训练、改进模型结构仍有继续提升空间。

GUI 部分将命令行实验结果转化为可交互展示系统，支持 demo 播放、波形与频谱可视化、自定义音频分离和指标总览。它使实验结果不再只停留在 CSV 和终端输出中，而是可以通过听觉和视觉方式直接验证分离效果。

后续改进方向包括：

\begin{enumerate}
    \item \textbf{扩大训练数据规模}：当前训练样本数量较小，模型对复杂说话人和噪声场景的泛化能力有限。增加更多说话人、更多语音片段和更多噪声类型有助于提升稳定性。
    \item \textbf{改进模型结构}：可以尝试 Conv-TasNet、Dual-Path RNN、SepFormer 等更强的语音分离模型，减少手工 STFT 特征依赖并提升端到端建模能力。
    \item \textbf{增强噪声鲁棒性}：Hard 实验中噪声显著影响分离质量，后续可加入不同 SNR 的噪声增强、混响增强或多条件训练。
    \item \textbf{完善 GUI 功能}：可继续加入批量分离、失败案例筛选、指标排序、音频下载按钮和模型对比页面，使 GUI 更接近完整实验分析平台。
    \item \textbf{补充主观评价}：除 SI-SDR 外，可以加入 PESQ、STOI 或人工听感评价，更全面地衡量分离语音的可懂度和自然度。
\end{enumerate}

总体而言，本实验实现了从数据生成、模型训练、结果评估到 GUI 展示的完整流程，能够较好地说明基础盲语音分离方法在简单和困难条件下的表现差异。
